\section{Institutional Interpretation and Limitations}
\label{sec:discussion}

Our extensive form is a stylized description of
\emph{pre-international technical coordination},
not a description of the voting rules of ISO or
IEC. It is plausible for a pair of governments
or nationally backed standards organizations to
coordinate an interface and later consider
broader participation. Historical motivations
include French--German cooperation during the
development of GSM before wider deployment
commitments, and collaboration between the
CHAdeMO Association and the China Electricity
Council on a next-generation charging interface
\citep{GSMAHistory,CHAdeMO2020}.
The latter also illustrates why an interoperable
physical interface need not fix every protocol
or residual product feature. Neither history
is evidence that the exact voting protocol in
Section~\ref{sec:model} operated, or that a
third country historically rejected accession
because of the design mechanism we model.

A positive $g$ in the model represents an
economically binding compatibility decision
that changes cross-product substitution and
network use in the ensuing product market.
Mere information sharing between technical
committees without deployed compatibility
would \emph{not} be sufficient to give a
soft memorandum the meaning of $g>0$.
Actual charging networks have two-sided
adoption and installed-base dynamics, and
actual telecommunications standardization
involves institutions and actors beyond
our three identical countries. These features
are deliberately absent from
\eqref{eq:utility}--\eqref{eq:welfare}.

A useful diagnostic comes from comparing the
government agenda with a distinct hypothetical
procedure. If the three countries were presented
with a \emph{single} unanimously accepted or
rejected $\IS$ offer at $g=.05$, with $\SW$
as the fallback and no intermediate $\SU$,
Proposition~\ref{prop:pareto} establishes that
each country would prefer accepting under
the specified market continuation. The
inaccessibility result is therefore tied to
the \emph{bilateral-first pathway and its
outsider opportunity}; it is not a theorem
that internationally compatible technology
necessarily fails in every bargaining
institution. We have not modelled the
legal authority, commitment mechanism,
transfer scheme, or enforcement that could
make this alternative offer available.
It should not be read as an implementable
policy proposal.

Other boundaries are equally important.
The results are computer-assisted existence
claims at an exact rational witness and
an unquantified local parameter neighborhood,
not a broad calibration to real firms.
Three countries are symmetric; each firm
serves all three identical national markets;
network expectations concern contemporaneous
quantities; demand is quadratic; and the
design game uses a real once-paid
quadratic cost. The certified nonlinear
design check changes only one substitution
mapping, not the demand foundation or
sequential government protocol.
All-g equilibrium/root assertions rely on
rational interval methods, while the separately
built Lean artifact proves four pointwise
payoff comparisons at \emph{rational design
proxies}; it does not formalize the whole
Nash and subgame-perfect equilibrium analysis.

\section{Conclusion}\label{sec:conclusion}

Evaluating a prospective international
technical agreement with firms' product
attributes held fixed can mispredict which
government coalitions are feasible. In a
three-country bilateral-first standards
consent game, we construct and certify a
case in which post-policy private
differentiation rules out every positive
depth of regional or international
agreement, even though a hypothetical
international standard still strictly
improves each country's welfare over
fragmentation. A forecast that fixes
firms' designs at their fragmented-market
choices instead predicts a feasible
international agreement. The qualitative
existence contrast persists under one
separately certified nonlinear residual
design mapping, although the depth of a
feasible frozen-design offer changes
substantially. The contribution is a
specific interaction between private
design responses and a sequential
participation institution, not a
general prescription for international
standardization.
